\Article{22}
Although we are certain \textit{a priori} that the series for $V,X,Y,Z$ converge, nothing can be determined in advance about the rate of convergence. If the locations of the magnetic forces were confined to a moderate region around the center of the Earth, or if the magnetic fluids were distributed within the Earth in a way equivalent to that case, the series would have to converge very rapidly. The farther those locations extend toward the surface, however, and the more irregular the distribution, the more we must be prepared for slow convergence.

In practical applications, absolute accuracy is unattainable. Only a degree of approximation suited to the circumstances is required. The slower the convergence, the more terms will have to be taken into account to achieve a given degree of accuracy.

Now $P'$ contains three terms and therefore requires knowledge of three coefficients $g^{1,0},g^{1,1},h^{1,1}$. $P''$ requires five coefficients, $P'''$ seven, $P''''$ nine, and so on. Thus, if we regard $P',P'',P'''$ and so on as quantities of the first, second, third order, and so on, it is clear that if the calculation is to extend through quantities of order $n$ inclusive, the values of $nn+2n$ \SourceLeaf{150}coefficients must be determined, for example 24 if we wish to go as far as the fourth order.

\expandafter\def\csname HistoricalNoteText18\endcsname{In such a method of determination, these circumstances would have the least adverse effect if the eight points were distributed quite symmetrically on the surface of the Earth, that is, if they coincided with the vertices of a cube inscribed in the sphere of the Earth, or at least came very close to such an arrangement.}
Every given value of $X,Y$ or $Z$, for given values of $u$ and $\lambda$, gives us an equation relating the coefficients. Thus, when all the elements of terrestrial magnetism are known at a place, they give three equations. If we could assume that only the terms up to the fourth order remain appreciable, then, in theory, complete observations at eight distinct points would suffice to determine all the necessary coefficients. But that assumption is hardly admissible, and the random errors affecting all observations, together with the neglect of higher-order terms, could seriously distort the results of elimination\HistoricalNote{18}. To reduce the harmful effect of these circumstances, we would have to use, as a basis, many more observed quantities than there are unknowns, from widely separated points in all parts of the Earth, and derive the unknown quantities from them by the method of least squares. Monotonous as it would be to carry out such a task, since the equations are all linear, the extraordinary amount of work arising from the great number of unknown quantities and equations might deter even the boldest person doing the calculations from undertaking the work in this form at present, especially since unreliable observations or errors of calculation could creep in and spoil the result entirely.

